\section{Conclusion}
\label{sec:conclusion}

A dataset of 5,460 CFD simulations, spanning 35 geometries, 12 Reynolds numbers, and 13 Richardson numbers, was generated for mixed convection in an eccentric horizontal annulus with a rotating inner cylinder, validated against experimental data, a qualitative rotation-effect trend from an independent numerical study, and a grid-convergence study following standard GCI procedure. The Reynolds and Richardson numbers used to define this sweep replace the Rayleigh number and rotation rate of earlier studies of this configuration, a reparameterization justified on physical grounds: for a single rotating boundary, the rotational contribution to the flow is fully captured by the Reynolds number, so sweeping rotation rate and viscosity independently would be redundant. All 5,460 simulations converged under an adaptive, iterate-until-stable criterion; a gradient-boosted tree ensemble was the most accurate of three candidate architectures for a scalar Nusselt number surrogate, achieving $R^2=0.9996$ on a held-out split of the training grid. A true out-of-sample test against eighteen new numerical simulations showed a substantially higher error of $7.89\%$. This error concentrates specifically at elevated Reynolds number, not at the Richardson-number-controlled transition between buoyancy- and rotation-dominated flow ($r=0.72$ for error against $Re$, versus $r=0.20$ against $Ri$, with the two largest errors occurring at opposite ends of the $Ri$ range): the twelve-level Reynolds number grid, log-spaced for even coverage in relative terms, leaves absolute gaps between training levels that grow roughly fifty-to-seventy-fold across the sweep, and this coarser absolute resolution coincides with the region where $Nu$ responds most steeply to $Re$, compounding into larger interpolation error with no change in the underlying flow physics. The model's own bootstrap-ensemble uncertainty estimate, well-calibrated on the training grid, was badly overconfident under this same genuine extrapolation. For the spatial temperature and velocity field, a per-node network again outperformed both a DeepONet and a graph neural network built on the true mesh connectivity, confirming the same architecture ranking found under an earlier, differently parameterized version of this study.

Four findings are offered as broader methodological cautions. First, same-grid holdout accuracy can substantially overstate accuracy under genuine extrapolation; an independent out-of-sample simulation should be used to check this directly, as done here. Second, an ensemble-based uncertainty estimate well-calibrated in-distribution should not be assumed well-calibrated outside it. Third, an initial set of twelve out-of-sample points selected without regard to their output value left an unintended gap in achieved Rayleigh number, understating the true extrapolation error ($5.92\%$ against the eventual $7.89\%$) and overstating uncertainty calibration ($41.7\%$ against $27.8\%$) purely because the points that would have exposed larger errors were never sampled; such points should be chosen, or checked after the fact, for coverage of the achieved output range, not only the input ranges. Fourth, a sampling grid designed for even coverage in relative (log) terms can still leave coarse absolute resolution where the underlying response is steepest; the Reynolds-number-specific error pattern identified above was traced to exactly this effect rather than to a change in convection regime, and future sampling designs for surrogate-modeling studies of this kind should consider absolute, not only relative, resolution when the target function's sensitivity varies strongly across the swept range.
